\paragraph{Result R53: the residue trace and the complete collision map.}
For the unchanged monic polynomial \(H(S)=\chi(S)-t\), the complete
SGA trace chapter constructs the perfect residue functional
\(\lambda_t[P]=[S^{q-1}]\operatorname{rem}_H P\) and the actual
diagonal coevaluation
\[
 \mathcal C_H(X,Y)=\frac{H(X)-H(Y)}{X-Y},\qquad
 \operatorname{Tr}_{E_t/R[t]}M_P=\lambda_t(\chi'P).
\]
The fraction is a polynomial, its multiplication image is \(\chi'\),
and the chapter proves the inverse pairing identities without
dividing by a discriminant. The original logarithmic source
relation maps to \([\chi'P]_\chi\) before this trace is taken.
Its complete periodic diagonal resolution retains both odd and
positive even degrees.

The new comparison (SC26)--(SC30) proves the precise specialization
of those degrees. Over \(T=\mathbb C[t]\), put
\(\mathscr E=T[S]/(\chi-t)\) and \(M=\mathscr E/\chi'\mathscr E\).
The isomorphism \(\mathscr E\simeq\mathbb C[S]\), with
\(t\mapsto\chi(S)\), makes
\(0\to\mathscr E\xrightarrow{\chi'}\mathscr E\to M\to0\)
a free \(T\)-resolution. For \(\kappa=\mathbb C[t]/(t-t_0)\),
\[
 \operatorname{Tor}_1^T(M,\kappa)
 \xrightarrow{\sim}\operatorname{ann}_{E_{t_0}}(\chi'),
 \qquad [v]_{\chi'}\longmapsto[w]_{\chi-t_0},
 \quad(\chi-t_0)v=\chi'w.
\]
The proof gives both inverse directions and the exact \(S\)-action.
Thus every positive even group of the family is zero, while the
positive even groups at a collision are these explicit Tor modules;
the odd groups are \(M\otimes_T\kappa\). Each length-\(m\) local
factor retains both \(m-1\)-dimensional modules and their nilpotent
actions. The Gamma calculation also evaluates the entire period
Gram determinant as
\[
 \det(\Pi^*\Pi)=(2\pi|u|)^q
  \exp\left(2\operatorname{Re}
          \frac{\operatorname{Tr}\Phi_t(A+t\mathcal R)}u\right),
\]
with the complete phase-containing determinant still given by R52.

\paragraph{Result R54: exact constituent curvature and neighboring volumes.}
For an actual invariant injection \(I:F\hookrightarrow E\),
\(0<p=\dim F<q\), put
\(\ell[P]=[S^{q-1}]\operatorname{rem}_\chi P\),
\(\mathcal R=e_0\ell\), \(L=\ell I\),
\(Y=\Pi I\), \(H=Y^*Y\),
\(N=1-YH^{-1}Y^*\), and \(z=N\Pi e_0\).
The complete proof (SC9)--(SC25) identifies the invariant image
with the unchanged polynomial ideal \((g)/(\chi)\), proves
\(L\ne0\) and \(z(t)\ne0\) for every parameter, and computes
\[
 NY'=-\frac t u zL,\qquad
 \partial_{\bar t}\partial_t\log\det H
    =\frac{|t|^2}{|u|^2}\|z\|^2LH^{-1}L^*.
\]
The normal derivative has rank one at every nonzero parameter;
the normal acceleration has rank one at zero. The curvature
vanishes precisely quadratically at zero and is positive elsewhere.
Its factorization through the original section \(r_Ne_0L\),
quotient \(j_E\), and period observation \(N\Pi\) retains the
source kernel and its supported zero.

For the literal polynomial frames
\(J_g=[g,Sg,\ldots,S^{p-1}g]\), its prefix \(J_-\), and
\(J_+=[J_g,e_0]\), put
\(h_J=\det(J^*\Pi^*\Pi J)\). The full frame map,
Schur-complement and cofactor proof in (SC31)--(SC37) gives
\[
 \partial_{\bar t}\partial_t\log h_g
       =\frac{|t|^2h_+h_-}{|u|^2h_g^2}.
\]
Each volume factors exactly as \(h_J=v_{N,J}b_{N,J}\), where
\(v_{N,J}=\det(J^*G_NJ)\) retains the original theta metric
and \(b_{N,J}\) is its specified period-comparison determinant.
The empty prefix has determinant one. In codimension one,
\(\det J_+=(-1)^{q-1}\), so \(h_+\) is the full explicitly
evaluated period determinant above. These are the exact maps and
weights joining this curvature to the neighboring-volume calculation.

\paragraph{Result R55: the retained Hochschild and critical-line comparisons.}
The complete Connes--Consani comparison proves the actual restricted
factorization
\[
 V\xrightarrow{j}j(V)\xrightarrow{\operatorname{Tr}}\mathscr B,
 \quad \operatorname{Tr}j=\Theta,
 \quad\mathcal E=\tfrac12\mathcal U\Theta,
 \quad (D^2-D)e^{-\pi u^2}=\phi_*.
\]
The half-density map \(\mathcal U F=u^{1/2}F\) is the specified
half-line isometry. The cochain isomorphism uses the complete
two-moment source and the Laplacian's continuous inverse. The
critical-line map \(\Gamma:Q\to Q_{\rm crit}\) has, on each
original finite packet, the exact kernel
\(\ker(\Gamma\sigma_Z)=E_{Z,\rm off}\), including every
generalized block. The critical recovery is the explicit jet
\[
 \Xi_\rho[f]=\sum_{j=0}^{m-1}
      \frac{i^{-j}}{j!}f^{(j)}(\gamma)\epsilon^j,
 \qquad \rho=\tfrac12+i\gamma.
\]
The full note proves the reciprocal multiplier argument, the
recovery composite and the supported image of every killed vector.

The raw circle zero mode is retained as a separate smooth coordinate;
the specified Gaussian average of the original theta seed proves
that this line is in the closed source-generated module. Its exact
local flat-seminorm estimate and the printed large-circle limit
correction remain with that proof. Finally the original-metric
Laplacian identities
\[
 \mathsf L_k^*G-G\mathsf L_k=A^*W_k-W_kA,
 \qquad G\mathsf L_k=W_kA-A^*GA
\]
and their complete numerical-range argument are retained with the
actual metric control parameter. Neither the trace construction,
the finite-kernel calculation, nor the period determinant supplies
an unproved uniform bound for that parameter.
